\section{Experimental Setup}
\label{sec:protocol}
\textbf{Data and evaluation protocol.}
Adaptation uses 40 ESC-50 classes: 256 recordings, 1,024 relation queries (512 pairs), and 512 rehearsal queries. Ten held-out classes provide 80 recordings, 320 queries, and 160 relation-development pairs, with disjoint original Freesound IDs. System evaluation uses 400 SpotSound queries, 6,649 Clotho-Moment queries, and the released 100-query/77-audio UnAV subset~\cite{munakata2025amr}. UnAV scores clip intervals to audio duration (unclipped mIoU: 73.39). SpotSound reports 997 UnAV queries; our evaluation uses the released 100-query subset. Development and public-benchmark scores were used during method development.

\textbf{Training settings.}
SFT and RelTwin initialize from SpotSound-A on Audio Flamingo~3. LoRA~\cite{hu2022lora} uses rank~8, $\alpha=16$, dropout~0.1, and 20.19M parameters. Both run 256 AdamW~\cite{loshchilov2019adamw} updates at learning rate $5\times10^{-6}$, zero weight decay, gradient clipping 1.0, $\tau=1$, $\beta=0.5$, and RelTwin $\lambda=1$. Each update accumulates one inverse-query pair and one rehearsal query. The matched comparison shares initialization, data order, code, environment, and optimizer updates. RelTwin additionally computes four gradient-free candidate scores followed by weighted gradient recomputation, versus SFT's two positive-answer backpropagations. Both include rehearsal.

\textbf{Inference.}
SFT and RelTwin are evaluated after 256 updates. All generators use 16\,kHz mono audio, bfloat16, and greedy one-beam decoding without sampling, capped at 128 new tokens. The post-development timing diagnostic uses the frozen matched checkpoints (Section~\ref{sec:timing}).

\textbf{Reference models.}
Table~\ref{tab:training} includes paired SFT and the official backbone. The backbone scores 58.42 mIoU versus published 57.9; RelTwin's three-seed difference against it is $+0.75$ points (audio-group CI $[-0.24,1.74]$).

\textbf{Result reporting and statistics.}
RelTwin point results use seed~0, selected by SpotSound mIoU from five runs, with its matched SFT control. Table~\ref{tab:training}(c) separately summarizes seeds 0/1/2 as mean $\pm$ sample SD. Training comparisons use 20,000 audio-group bootstrap resamples for 95\% CIs. Relation comparisons resample 39 connected source-sharing groups (largest: 18 audios). CIs quantify audio-sample uncertainty conditional on the evaluated checkpoints; model selection is fixed across resamples.

\subsection{NOVA system configurations}
\label{sec:system}
NOVA (Necessity-Oriented Verification over Audio) combines candidate generation, verification, and refinement. SpotSound uses RelTwin data with Jensen--Shannon consistency between $p_+$ and label-swapped $p_-$ added to Eq.~\ref{eq:loss}. Setwise preference optimization (SetPO) continues for 64 updates at $2\times10^{-6}$ over six candidates: both windows, their expansions, enclosing span, and union. Soft targets combine set-IoU and interval-coverage quality. The RelTwin training comparison in Table~\ref{tab:training} uses the three-term objective in Eq.~\ref{eq:loss}.

Official, SFT, and SetPO seed-1 SpotSound predictions enter a LongNeedle-ESC50-fitted ridge router, using geometry and frozen query-existence scores after keeping/removing candidate intervals. Clotho and UnAV apply refinement to official predictions, whose initial mIoUs are 86.8543 and 73.4513. Refinement uses frozen SetPO/SpanTool features and ridge edit utilities: dynamic programming selects compatible boundaries, accepted when the 200-model bootstrap 5th percentile exceeds 0.02. Event count is preserved; radius $\min(0.25\,\mathrm{s},0.625\,\text{event duration})$ is shared. AudioGrounding/AEGBench development evaluations selected the duration coefficient before these applications. Table~\ref{tab:main} reports these NOVA configurations.
